\color{unidarkgreen}{\subsection{Контроль синхронизма и напряжений}}
\color{black}

\subsubsection{Общие сведения}

\begin{enumerate}[label=\arabic{section}.\arabic{subsection}.\arabic{subsubsection}.\arabic*,  align=left, labelsep=4pt, leftmargin=0pt, itemindent=57pt]

\item
Функциональный блок контроля синхронизма и напряжений (КСН) предназначен для
формирования разрешающих сигналов, соответствующих условиям включения, для АПВ, оперативного
включения.

\item
Основные режимы включения и контролируемые для них условия приведены в таблице \ref{}.
Режимы включения с контролем синхронизма (КС) и с улавливанием синхронизма
(УС) вводятся соответствующими программными ключами \textbf{«КС В1(ЛВ) ВН»} и \textbf{«УС В1(ЛВ) ВН»}.

\item
Если разность частот напряжений шин и трансформатор не превышает уставку \textbf{«df макс КС»},
то разрешается включение с контролем синхронизма. Для режима включения с КС также контролируется:
\begin{itemize}
\item наличие напряжения на обоих энергообъектах (шины и трансформатор);
\item разность напряжений шин и трансформатора, не превышающая по модулю значения заданной уставки;
\item разность фаз напряжений шин и трансформатора, не превышающая по модулю значения заданной уставки.
\end{itemize}

\item
При выводе программного ключа \textbf{«КС В1(ЛВ) ВН»} включение разрешается без контроля частоты,
а только по уровню напряжений и разности их фаз.

\item
Если разность частот превышает уставку \textbf{«df макс КС»}, но не выходит за пределы уставки \textbf{«df
пред. доп.»}, то разрешается включение с улавливанием синхронизма. Сигнал включения с
улавливанием синхронизма формируется с опережением на время включения выключателя. Режим
включения с УС можно вывести программным ключом \textbf{«УС В1(ЛВ) ВН»}.

\item
Наличие или отсутствие напряжения на трансформатор контролируется измерительными
органами максимального и минимального напряжения, пороги срабатывания которых задаются
уставками \textbf{«Uнал Т В1(ЛВ) ВН»} и \textbf{«Uотс Т В1(ЛВ) ВН»} соответственно. Пороги срабатывания
измерительных органов для контроля напряжения на энергообъекте Ш (шины) задаются уставками
\textbf{«Uнал Ш В1(ЛВ) ВН»} и \textbf{«Uотс Ш В1(ЛВ) ВН»}.

\small
\begin{longtable}{|>{\centering\arraybackslash}m{5.3cm}|>{\centering\arraybackslash}m{11.6cm}|}
% Первая страница: полный заголовок
\caption{Условия включения\hfill\vspace{-0.5\baselineskip}}\label{desc:tbl1}\\
\hline
\rowcolor{gray!30}
Режим включения & Условия включения \\
\hline
\endfirsthead

% Промежуточные страницы: «Продолжение таблицы»
\caption*{\hspace{3pt}\emph{Продолжение таблицы \ref{desc:tbl1}}\hfill\ }\vspace{-0.5\baselineskip} \\
\hline
\rowcolor{gray!30}
Режим включения & Условия включения \\
\hline
\endhead

ОНТ+ННШ &   \vspace{5pt}
			\begin{itemize}
			\item $ \textrm{U}_\textrm{T} $ меньше уставки \textbf{«Uотс Т В1(ЛВ) ВН»};
			\item $ \textrm{U}_\textrm{Ш} $ больше уставки \textbf{«Uнал Ш В1(ЛВ) ВН»}.
			\end{itemize} \\
			\vspace{10pt}
\hline
ННШ +ННТ & 
			\begin{itemize}
			\item $ \textrm{U}_\textrm{T} $ больше уставки \textbf{«Uнал Т В1(ЛВ) ВН»};
			\item $ \textrm{U}_\textrm{Ш} $ меньше уставки \textbf{«Uотс Ш В1(ЛВ) ВН»}.
			\end{itemize} \\
\hline
КС(УС)	&
			1) Включение контролем синхронизма: 
			\begin{itemize}
			\item $ \textrm{U}_\textrm{T} $ больше уставки \textbf{«Uнал Т В1(ЛВ) ВН»};
			\item $ \textrm{U}_\textrm{Ш} $ больше уставки \textbf{«Uнал Ш В1(ЛВ) ВН»};
			\item разность напряжений трансформатора и шин не превышает по модулю значения заданной уставки \textbf{«dU макс КС(УС)»};
			\item разность фаз напряжений трансформатора и шин не превышает по модулю значения заданной уставки \textbf{«dФ макс КС»};
			\item разность частот напряжений трансформатора и шин не превышает значения заданной уставки \textbf{«df макс КС»};
			\end{itemize}
			
			2) Включение с улавливанием синхронизма: 
			\begin{itemize}
			\item $ \textrm{U}_\textrm{T} $ больше уставки \textbf{«Uнал Т В1(ЛВ) ВН»};
			\item $ \textrm{U}_\textrm{Ш} $ больше уставки \textbf{«Uнал Ш В1(ЛВ) ВН»};
			\item разность напряжений трансформатор и шин не превышает по модулю значения заданной уставки \textbf{«dU макс КС(УС)»};
			\item разность частот напряжений трансформатора и шин превышает значение уставки \textbf{«df макс КС»};
			\item разность частот напряжений трансформатора и шин не превышает значения заданной уставки \textbf{«df пред. доп.»};
			\item за время \textbf{«Т опер. вкл. В1(ЛВ) ВН»} с момента подачи сигнала на включение выключателя расхождение углов между векторами напряжений на трансформаторе и шинне превысит 10º.
			\end{itemize} \\
\hline

\end{longtable}
\normalsize	

\begin{figure}[H]
\centering
\includegraphics[width=0.8\textwidth,height=0.8\textheight,keepaspectratio]{\fbpath/761.3431 БКУ/img/BKU.pdf}
\captionof{figure}{Функциональная схема команд управления}\label{pic:BKU}
\end{figure}


\small
\begin{longtable}{|>{\centering\arraybackslash}m{5.3cm}|>{\centering\arraybackslash}m{4.2cm}|>{\centering\arraybackslash}m{1cm}|>{\centering\arraybackslash}m{3.3cm}|>{\centering\arraybackslash}m{1.8cm}|}
% Первая страница: полный заголовок
\caption{Перечень уставочных параметров команд управления\hfill\vspace{-0.5\baselineskip}}\label{desc:tbl1}\\
\hline
\rowcolor{gray!30}
Наименование уставки & Диапазон & Шаг & Значение по умолчанию & Единица измерения \\
\hline
\endfirsthead

% Промежуточные страницы: «Продолжение таблицы»
\caption*{\hspace{3pt}\emph{Продолжение таблицы \ref{desc:tbl1}}\hfill\ }\vspace{-0.5\baselineskip} \\
\hline
\rowcolor{gray!30}
Наименование уставки & Диапазон & Шаг & Значение по умолчанию & Единица измерения \\
\hline
\endhead

%===m>ДТЗ|-
%===m
\end{longtable}
\normalsize	

\end{enumerate}